\section{Finite choice}

\import{set/finite.tex}

Choice over a finite domain follows by adjoining one selected value at a
time.  This result does not use the global choice operator.

\begin{theorem}\label{finite_choice}
    Suppose $A$ is finite and $R$ is a relation.
    Suppose for every $x\in A$ there exists $y$ such that
        $x\mathrel{R}y$.
    Then there exists $f$ such that $f$ is a function on $A$ and
        for every $x\in A$ we have $x\mathrel{R}f(x)$.
\end{theorem}
\begin{proof}
    Let $C=\{B\in\pow{A}\mid
        \text{there exists $f\in\funs{B}{\ran{R}}$ such that
            for every $x\in B$ we have $x\mathrel{R}f(x)$}\}$.

    $\emptyset\in\pow{A}$ by \obv{powerset_bottom}.
    $\emptyset\in\funs{\emptyset}{\ran{R}}$ by
        \obv{emptyset_is_function_on_emptyset,
            codom_of_emptyset_can_be_anything,funs_intro}.
    For every $x\in\emptyset$ we have
        $x\mathrel{R}\apply{\emptyset}{x}$ by \auto{emptyset}.
    Hence $\emptyset\in C$.

    Show that for all $B,b$ such that $B\in C$ and $b\in A$ we have
        $\cons{b}{B}\in C$.
    \begin{subproof}
        Fix $B,b$.
        Assume $B\in C$ and $b\in A$.
        $B\subseteq A$ by \obv{powerset_elim}.
        Hence $\cons{b}{B}\subseteq A$ by
            \obv{cons_subseteq_intro}.
        Thus $\cons{b}{B}\in\pow{A}$ by \obv{powerset_intro}.
        \begin{byCase}
        \caseOf{$b\in B$.}
            $\cons{b}{B}=B$ by \obv{cons_absorb}.
        \caseOf{$b\notin B$.}
            Take $f\in\funs{B}{\ran{R}}$ such that for every $x\in B$
                we have $x\mathrel{R}f(x)$.
            Take $y$ such that $b\mathrel{R}y$.
            Hence $(b,y)\in R$.
            Thus $y\in\ran{R}$ by \obv{ran_intro}.
            Let $g=\cons{(b,y)}{f}$.
            $g\in\funs{\cons{b}{B}}{\ran{R}}$ by
                \auto{extend_function_at_new_point}.
            Show that for every $x\in\cons{b}{B}$ we have
                $x\mathrel{R}g(x)$.
            \begin{subproof}
                Fix $x\in\cons{b}{B}$.
                $x=b$ or $x\in B$ by \obv{cons_iff}.
                \begin{byCase}
                \caseOf{$x=b$.}
                    $(b,y)\in g$ by \auto{cons_iff}.
                    $g$ is a function by \obv{funs_is_function}.
                    Hence $g(b)=y$ by \obv{function_apply_intro}.
                \caseOf{$x\neq b$.}
                    Hence $x\in B$.
                    $f$ is a function and $\dom{f}=B$ by
                        \obv{funs_is_function,funs}.
                    Hence $x\in\dom{f}$.
                    Thus $(x,f(x))\in f$ by
                        \obv{function_apply_elim}.
                    Hence $(x,f(x))\in g$ by \auto{cons_iff}.
                    $g$ is a function by \obv{funs_is_function}.
                    Hence $g(x)=f(x)$ by
                        \obv{function_apply_intro}.
                \end{byCase}
            \end{subproof}
        \end{byCase}
    \end{subproof}

    Hence $A\in C$ by \obv{finite_induction}.
    Take $f\in\funs{A}{\ran{R}}$ such that for every $x\in A$ we have
        $x\mathrel{R}f(x)$.
    $f$ is a function on $A$ by \obv{funs_is_function,funs}.
    Follows by \auto.
\end{proof}
